\section*{第 \ref{ch:g5:food-webs} 章——食物网}

\begin{solution}{exo:g5:food-webs:1}
把一个\omterm{def:g2:where-animals-live:habitat}{栖息地}的全部\omterm{def:g3:food-chains:chain}{食物链}一次画在一起形成的网络——
每个\omterm{def:g5:biodiversity:species}{物种}只画一次，每一种取食关系用一个箭头表示。\omterm{def:g3:food-chains:chain}{食物链}
之所以相互交叉，是因为大多数\omterm{def:g1:animals-around-us:animal}{动物}既吃多种\omterm{def:g1:living-or-not:living}{生物}，
也被多种\omterm{def:g1:living-or-not:living}{生物}吃掉。
\end{solution}

\begin{solution}{exo:g5:food-webs:2}
\omterm{def:g5:food-webs:roles}{生产者}：\omterm{def:g1:plants-around-us:plant}{植物}（草、三叶草、橡树叶）。\omterm{def:g5:food-webs:roles}{消费者}：
\omterm{def:g1:animals-around-us:animal}{动物}——\omterm{def:g2:what-animals-eat:kinds}{植食性动物}如兔子，\omterm{def:g2:what-animals-eat:kinds}{肉食性动物}如鹰。\omterm{def:g5:food-webs:roles}{分解者}：
以死亡遗体为食的腐食\omterm{def:g1:living-or-not:living}{生物}（图中未画出——
它们在整张网的
影子下默默工作）。
\end{solution}

\begin{solution}{exo:g5:food-webs:3}
蓝山雀以\omterm{ex:g2:animals-grow-up:butterfly}{毛毛虫}和蚱蜢为食，
它又被雀鹰吃掉。
\end{solution}

\begin{solution}{exo:g5:food-webs:4}
草 $\rightarrow$ 蚱蜢 $\rightarrow$ 蜥蜴
$\rightarrow$ 鹰，以及草 $\rightarrow$ 蚱蜢
$\rightarrow$ 蓝山雀 $\rightarrow$ 雀鹰。
\end{solution}

\begin{solution}{exo:g5:food-webs:5}
\omterm{def:g5:food-webs:roles}{生产者}很多，\omterm{def:g2:what-animals-eat:kinds}{植食性动物}较少，\omterm{def:g2:what-animals-eat:kinds}{肉食性动物}更少。
因为每一层都以下面一层为食：取食者的数量
永远不可能超过可供它们吃的食物。
\end{solution}

\begin{solution}{exo:g5:food-webs:6}
因为一对鹰需要一大片草地所能提供的
蚱蜢和足够多的蜥蜴作为食物——金字塔的顶端
立足于下方宽阔的各层之上。
\end{solution}

\begin{solution}{exo:g5:food-webs:7}
它们的捕食者——狐狸、鹰——吃得更好，数量随之增多。
额外的捕猎使兔子的过剩数量熔化，
种群数量重新趋于平衡。
\end{solution}

\begin{solution}{exo:g5:food-webs:8}
开放性答案。一个可以成立的花园\omterm{def:g5:food-webs:web}{食物网}：生菜和草
（\omterm{def:g5:food-webs:roles}{生产者}）$\rightarrow$ 蜗牛、蚜虫、麻雀；
蜗牛 $\rightarrow$ 刺猬；蚜虫 $\rightarrow$ 瓢虫；
麻雀和瓢虫 $\rightarrow$ 邻居家的猫或
雀鹰。
\end{solution}

\begin{solution}{exo:g5:food-webs:9}
第一步：蚱蜢因被吃得少而增加；鹰
因少了一种猎物，更多地依赖兔子（它的备选
箭头）。第二步：更多的蚱蜢压低草的数量，
更多猎捕兔子使兔子数量下降——从而缓解
草的压力。\omterm{def:g5:food-webs:web}{食物网}随之调整，部分吸收了这一
损失；没有备选箭头的\omterm{def:g5:biodiversity:species}{物种}本会受到
最大伤害，而这里每个\omterm{def:g5:biodiversity:species}{物种}都有备选。
\end{solution}

\begin{solution}{exo:g5:food-webs:10}
清除鹰类之后，老鼠和田鼠失去主要捕食者，
数量增加；随后激增的啮齿\omterm{def:g1:animals-around-us:animal}{动物}吃掉了
收获的粮食——比鹰类曾经造成的损失还多。
倡导者只预见到第一步就停了手；\omterm{def:g5:food-webs:web}{食物网}实际上运行了
两步。
\end{solution}

\begin{solution}{exo:g5:food-webs:11}
橡树的损失自下而上冲击依赖它的\omterm{def:g1:living-or-not:living}{生物}：\omterm{ex:g2:animals-grow-up:butterfly}{毛毛虫}
（橡树叶是它唯一的箭头）挨饿，随后蓝山雀
完全依赖蚱蜢，随后雀鹰也感受到
数量减少——这是一场自下而上的“干旱”，就像干燥的夏季那样。
雀鹰的损失则自上而下传递：蓝山雀
因被捕食减少而增加，压低\omterm{ex:g2:animals-grow-up:butterfly}{毛毛虫}和蚱蜢的数量
——这就是清除捕食者的故事。这里鹰的损失
更容易被吸收：它的猎物还有其他捕食者，
扰动逐层减弱，而橡树的
依赖者完全没有备选箭头——\omterm{ex:g2:animals-grow-up:butterfly}{毛毛虫}的
整个层级只依靠一种\omterm{def:g5:food-webs:roles}{生产者}。
\end{solution}
